% =====================================================================
% Ensaio 1 — Seção 6: Diagnóstico de poder e limite superior
%
% ⚠️ ESTA SEÇÃO NÃO É "RESULTADOS", e o título não é eufemismo.
%
% O alvo declarado da qualificação é documento de PROJETO: identificação
% fechada, dose verificada, ameaças mapeadas, NÃO resultado estimado. A
% estimação foi levada até o fim, e o que ela mostra é que a análise de poder
% ex ante estava certa — o desenho não distingue o efeito procurado do zero.
% Isso é informação sobre o DESENHO, e é nessa chave que entra aqui.
%
% ⚠️ A pré-especificação (docs/pre-especificacao.md, fechada em 2026-09-21,
% commit c386514) declara em §5: se o desenho não falsificar, o que se escreve é
% LIMITE SUPERIOR INFORMATIVO. Nenhum número desta seção é apresentado como
% achado, e a §6 daquele documento declara o que é confirmatório e o que é
% exploratório. Esta seção respeita as duas coisas.
% =====================================================================

\section{Diagnóstico de poder e limite superior}
\label{sec:diagnostico}

\subsection{O que foi fixado antes de olhar}
\label{sub:prespec-fechada}

A análise que segue foi declarada antes de qualquer estimativa, em documento
datado e versionado. O compromisso relevante é o critério de falsificação:
fixou-se o piso de efeito esperado em \textbf{15 g} e declarou-se que, se a
meia-largura do intervalo de confiança não ficasse abaixo desse piso, o que se
escreveria seria \textbf{limite superior informativo}, não achado.

\paragraph{De onde vem esse piso.}
A âncora é \citet{reynier2025glyphosate}, que estimam redução de \textbf{23 a
32 g} no peso ao nascer à intensidade média de exposição ao glifosato no meio
rural norte-americano, com estimativa central de 29,8 g. O piso adotado aqui é
menor que o extremo inferior daquela faixa, e a diferença é deliberada: supõe-se
atenuação, porque proibir o \emph{método} aéreo não remove a molécula quando o
produtor substitui o avião pelo trator. Piso menor torna a falsificação mais
difícil, não mais fácil, e essa é a razão de tê-lo escolhido.

O transporte da âncora tem limite declarado. Aqueles autores medem o efeito
de \emph{acrescentar} glifosato; aqui mede-se o de \emph{remover} um método de
aplicação, e de outra classe química --- procimidona e carbaril, não glifosato
(§\ref{sec:background}). A âncora é analógica, e é sob essa ressalva que entra.
\citet{dias2023down}, o template de desenho deste trabalho, não serve de âncora
de magnitude: o resultado principal deles é de mortalidade infantil, não de peso
em gramas.

\paragraph{Uma ressalva de integridade que o documento faz de si mesmo.}
Este não é pré-registro no sentido estrito, e afirmá-lo seria falso. A janela em
que nenhuma fonte real havia sido tocada fechou antes da assinatura, e o
histórico de versionamento prova a data. O que existe é um \textbf{plano de
análise escrito depois dos diagnósticos de desenho e antes de qualquer
estimação}: conheciam-se a distribuição de dose, o desvio-padrão das tendências
municipais, o efeito mínimo detectável e o primeiro estágio do instrumento ---
tudo pré-período ou aritmética de desenho --- e não se conhecia nenhum
coeficiente. A fronteira que o documento protege é entre \emph{desenho} e
\emph{efeito}, não entre nenhum dado e dado, e essa fronteira é verificável.

\subsection{A curva dose-resposta}
\label{sub:curva}

O estimador é o de \citet{callaway2024continuous} com \emph{sieve}
não-paramétrico, sobre o painel colapsado em pré e pós --- colapso que o próprio
estimador exige, e que descarta as coortes cuja gestação atravessa o ban
(9,4\% das células). O grupo de comparação é o \(d = 0\) pela construção que
usa baixa aptidão agroclimática, e a definição alternativa entra como
sensibilidade.

\paragraph{O que se reporta aqui não é o alvo que a pré-especificação declarou
primário, e a troca está registrada.}
Aquele documento fixou o \(ACR(d)\) --- a curva --- como parâmetro principal, e
designou o \(ATT(d\,|\,d)\) como \emph{sensibilidade}. O que segue inverte os
dois. A razão é o parágrafo sobre a não-interpretabilidade da derivada agregada,
mais abaixo nesta mesma subseção, e a inversão entra na tabela de desvios da §8
da pré-especificação com data de 22 de setembro de 2026. Dizer isto antes do
número, e não depois, é o que impede que a troca pareça conveniência.

O \(ATT(d\,|\,d)\) agregado é de \textbf{\(-36{,}19\) g}, com erro-padrão de
15,47 g e intervalo de confiança de 95\% em \([-66{,}51;\ -5{,}87]\). A
meia-largura é de \textbf{30,32 g}, contra o piso pré-especificado de 15 g.

\begin{quote}
\textbf{Pelo critério fixado antes de olhar, isto não é achado.} É limite
superior informativo: o desenho é compatível com efeito nulo e com efeito de
magnitude compatível com a literatura, e não distingue entre os dois.
\end{quote}

\paragraph{A sensibilidade ao grupo de comparação é estreita.}
Trocar a construção do \(d = 0\) desloca o agregado de \(-36{,}19\) para
\(-35{,}33\) g. O \emph{sieve} centra a curva na variação média do grupo de dose
zero, de modo que contaminação desse grupo deslocaria o nível inteiro; aqui o
deslocamento é inferior a 1 g, porque apenas um município muda de lado.

\paragraph{O \(ACR(d)\) agregado não é interpretável, e isso é achado
metodológico.}
A derivada da curva é fortemente não linear e inverte de sinal ao longo da dose.
Na faixa de dose inferior a 0{,}1\% --- onde estão 48 dos 169 pontos --- ela é
da ordem de \(-9.100\) g por unidade de dose, com erro-padrão \emph{maior} que a
estimativa. O agregado que o pacote reporta é média simples sobre os pontos, e
portanto é dominado justamente pela região onde a derivada não é identificada.
A magnitude, ademais, é artefato de unidade: naquela faixa a dose varia cerca de
0{,}0009, de modo que \(-9.100\) g \emph{por unidade de dose} implica algo como
\(-8{,}6\) g de efeito. Nenhum dos 169 pontos exclui o zero, em nenhuma faixa. O
que se reporta, portanto, é o \(ATT(d\,|\,d)\), que é estável entre faixas.

\subsection{Inferência com poucos clusters tratados}
\label{sub:inferencia-poucos}

O \emph{bootstrap} multiplicador do estimador principal é assintótico no número
de clusters, e com dezessete municípios de dose alta esse intervalo é estreito
demais para ser levado a sério. A Tabela~\ref{tab:inferencia} põe três
procedimentos lado a lado, sobre a mesma forma reduzida.

\begin{table}[htbp]
\centering
\caption{Três procedimentos de inferência sobre a mesma forma reduzida}
\label{tab:inferencia}
\small
\begin{tabular}{lrl}
\toprule
\textbf{Procedimento} & \textbf{\(p\)} & \textbf{Supõe} \\
\midrule
Cluster-robusto           & 0,736 & assintótico no nº de clusters \\
\emph{Wild cluster bootstrap} & 0,757 & \citet{cameron2008wild}; sub-rejeita com poucos tratados \citep{canay2021wild} \\
Inferência por aleatorização & 0,793 & nenhum assintótico \\
\bottomrule
\end{tabular}
\end{table}

A lógica é a de \citet{conley2011}: não confiar no assintótico quando as
unidades tratadas são poucas. A adaptação é fiel na lógica e não no
procedimento --- Conley e Taber tratam tratamento binário com poucas mudanças de
política, e aqui o tratamento é contínuo, de modo que a versão exata é permutar
a dose. O intervalo de 95\% por inversão é \([-46{,}45;\ +57{,}85]\).

\textbf{O informativo é a concordância.} Este arranjo existe para expor
distância entre procedimentos: quando o ingênuo rejeita e a aleatorização não, o
resultado não é ``significante'', é ``o desenho não distingue''. Aqui não há
distância a expor, porque não há precisão falsa a desmentir --- nenhum dos três
rejeita, e o coeficiente da forma reduzida (\(+6{,}81\)) está abaixo do efeito
mínimo detectável do próprio desenho (33{,}4 g).

\paragraph{A família confirmatória, corrigida.}
A pré-especificação declara duas hipóteses confirmatórias --- peso médio e óbito
fetal --- e manda corrigir a família por Holm, que controla a taxa de erro por
família sem a perda de potência do Bonferroni. Corrigidos, os \(p\) do peso vão
de 0,394 para 0,427 no teste unilateral que a direção declarada pede, e de 0,794
para 0,866 no bilateral; os do óbito fetal, de 0,214 para 0,427 e de 0,433 para
0,866. \textbf{Nenhuma hipótese rejeita a 5\%, antes ou depois da correção.}

A correção não muda a conclusão, e é precisamente por isso que vale registrá-la:
o compromisso era de procedimento, e um procedimento que só se cumpre quando o
resultado não depende dele não é procedimento. A pré-especificação não declara
qual dos três estimadores de incerteza alimenta o Holm --- lacuna dela, não
escolha deste texto ---, de modo que os três são reportados.

\subsection{O placebo tem poder? Quanto?}
\label{sub:pretrends}

A credibilidade do desenho repousa sobre o teste de tendências paralelas, e um
teste que ``passa'' pode estar passando por falta de poder, não por ausência de
violação \citep{roth2023trending}. Medir esse poder deixa de ser refinamento e
vira pré-requisito.

\paragraph{Uma conta ingênua que engana.}
A meia-largura mediana dos \emph{leads} individuais é de cerca de 86 g, contra o
piso de 15 g --- razão de quase seis para um, que sugeriria um teste inútil. A
conta está errada: o teste de tendência usa os 46 \emph{leads}
\textbf{conjuntamente} para detectar inclinação linear, e o teste conjunto é
muito mais potente que qualquer \emph{lead} isolado.

Feita a conta certa, uma tendência linear de \textbf{0,238 g por mês} é
detectada em 80\% das vezes; a 90\%, exige 0,429 g por mês. Traduzidas para a
janela de quatro anos de pós-tratamento, correspondem a \textbf{11,4 g} e
\textbf{20,6 g} de viés acumulado.

\begin{quote}
A frase honesta, portanto, tem o valor preenchido: \textbf{o placebo exclui
tendências confundidoras capazes de fabricar mais de cerca de 11 g de efeito
espúrio, e não exclui abaixo disso.} Ele não é cego, e também não é
tranquilizador --- 11 g é da ordem do efeito procurado.
\end{quote}

\paragraph{Uma limitação do estimador que restringe o que se pode dizer.}
Com coorte única --- que é o caso de um ban simultâneo --- o estimador principal
devolve estimativas de \emph{event study} apenas para os períodos anteriores ao
tratamento; os posteriores saem indefinidos. O que existe é, portanto, um teste
de pré-tendências, e não um \emph{event study}: a trajetória dinâmica pós-ban
não é recuperável por esta via.

\subsection{O preço da hipótese de paralelismo}
\label{sub:honestdid}

\citet{rambachan2023credible} respondem à pergunta seguinte: admitida violação
de paralelismo de tamanho \(\bar{M}\) relativo à maior violação pré observada,
que intervalo sobrevive? Como o estimador principal não entrega períodos pós, a
aplicação usa um \emph{event study} de efeitos fixos com tratamento
\textbf{binário} --- decil superior contra o grupo \(d = 0\) ---, agregado a
ano. É outro estimando, e o resultado abaixo é sobre ele.

\begin{table}[htbp]
\centering
\caption{Sensibilidade a violações de tendências paralelas}
\label{tab:honestdid}
\small
\begin{tabular}{lrr}
\toprule
\(\bar{M}\) & \textbf{IC inferior} & \textbf{IC superior} \\
\midrule
0{,}00 (paralelismo perfeito) & \(-72{,}5\) & \(+24{,}5\) \\
0{,}50 & \(-130{,}4\) & \(+79{,}5\) \\
1{,}00 & \(-217{,}4\) & \(+161{,}4\) \\
2{,}00 & \(-395{,}3\) & \(+339{,}3\) \\
\bottomrule
\end{tabular}
\end{table}

\textbf{O valor de quebra é zero}, porque o intervalo já cobre o zero sob
paralelismo perfeito. A leitura correta não é que o resultado seja frágil: é que
não havia significância a ser destruída. O procedimento não conserta poder ---
ele precifica a hipótese, e o preço aqui é que a largura quase quadruplica
quando se admite violação do tamanho da maior violação pré observada.

\paragraph{Uma observação sobre o uso de efeitos fixos de dois sentidos.}
A literatura recente adverte contra o \emph{event study} de efeitos fixos sob
adoção escalonada \citep{goodmanbacon2021,sunabraham2021}. Essa advertência
\textbf{não se aplica a este uso}: o ban é simultâneo e existe grupo
nunca-tratado limpo, de modo que não há a comparação entre unidades já tratadas
que produz o viés. É o caso em que o estimador é apropriado.

\subsection{A hipótese de que o alvo primário precisava, e que não havia sido
sondada}
\label{sub:spt}

As duas sensibilidades anteriores operam sobre o \emph{event study} binário, em
níveis, e falam portanto do paralelismo tradicional. O alvo que a
pré-especificação declarou principal --- a curva --- precisa de outra hipótese, o
\emph{strong parallel trends}, e essa não havia sido sondada por nenhum
procedimento deste trabalho.

\paragraph{A afirmação de que ela não é testável era forte demais.}
O material de projeto registrava que nenhum \emph{event study} de \emph{leads}
poderia falar dessa hipótese, porque ela envolve trajetórias sob doses não
recebidas, e no pré-período ninguém é tratado. O raciocínio está quase todo
certo, e a conclusão não. \citet{callaway2024continuous} observam, na §6.3, que
antes do tratamento as duas versões do paralelismo têm a \emph{mesma}
implicação --- a relação entre a variação do desfecho e a dose deve ser nula ---
e computam exatamente essa verificação na própria aplicação, onde ela
\textbf{rejeita}. A formulação correta é mais estreita: o placebo não
\emph{isola} a hipótese forte, porque passar é compatível com as duas; mas pode
\textbf{falsificá-la}.

\paragraph{O teste, e o que ele devolve.}
Finge-se que o banimento ocorreu em data do pré-período e roda-se a mesma forma
reduzida, com a mesma distância entre janelas do desenho real. A janela de
2015--2018 comporta três cortes com doze meses de cada lado.

\begin{table}[htbp]
\centering
\caption{Inclinação em dose antes do banimento: três cortes placebo}
\label{tab:spt}
\small
\begin{tabular}{lrrrr}
\toprule
& \multicolumn{2}{c}{\textbf{Peso médio}} & \multicolumn{2}{c}{\textbf{Óbito fetal}} \\
\cmidrule(lr){2-3}\cmidrule(lr){4-5}
\textbf{Corte placebo} & \textbf{incl.} & \textbf{\(p\)} & \textbf{incl.} & \textbf{\(p\)} \\
\midrule
2016-01 / 2016-11 & 25,76 & 0,371 & 0,00086 & 0,864 \\
2016-07 / 2017-05 & 31,07 & 0,239 & 0,00074 & 0,852 \\
2017-01 / 2017-11 & 12,96 & 0,608 & 0,00021 & 0,959 \\
\midrule
\emph{desenho real} & \emph{6,81} & \emph{0,794} & \emph{--0,00244} & \emph{0,433} \\
\bottomrule
\end{tabular}
\end{table}

Nenhum dos três rejeita, em nenhum dos dois desfechos. A implicação comum não é
violada onde se consegue olhar, e isso é o mais forte que este teste consegue
ser --- ele não valida a hipótese.

\begin{quote}
\textbf{E a coluna do peso traz o diagnóstico mais duro desta seção.} Os três
cortes placebo produzem inclinação \emph{maior em módulo} que a do desenho real:
25,8, 31,1 e 13,0 contra 6,8. O ruído de pré-período do próprio painel é maior
que o efeito estimado. Não é que o teste não rejeite por o desenho ser limpo; é
que ele não rejeitaria de todo modo.
\end{quote}

\paragraph{E a assimetria entre os dois desfechos é informativa.}
No óbito fetal a relação se inverte: nenhum dos três placebos alcança a
inclinação do desenho real. Onde o peso não separa sinal de ruído de
pré-período, o óbito fetal separa --- o que reforça que ele entrou como segundo
confirmatório por razão de desenho, e não por acessório.

\subsection{Controle sintético, e por que o resultado que cruza a linha não é
reportado como achado}
\label{sub:sdid}

O desenho previa controle sintético com diferenças-em-diferenças
\citep{arkhangelsky2021sdid} para o agregado estadual --- Ceará contra um Ceará
sintético construído com outros estados. Essa versão não é executável com os
dados adquiridos, que cobrem apenas municípios cearenses. O que se roda é a
versão intraestadual: municípios de dose alta contra sintético construído com os
de dose zero. É outra pergunta, e a ameaça também é outra --- se o grupo \(d=0\)
estiver contaminado por dispersão aérea para controle vetorial, o sintético é
construído com unidades parcialmente tratadas, e o método herda o problema em
vez de corrigi-lo.

\begin{table}[htbp]
\centering
\caption{Três estimadores sobre a mesma matriz de painel}
\label{tab:sdid}
\small
\begin{tabular}{lrrr}
\toprule
\textbf{Estimador} & \textbf{\(\tau\) (g)} & \textbf{EP} & \textbf{IC 95\%} \\
\midrule
Synthetic DiD       & \(-20{,}16\) & 18{,}50 & \([-56{,}43;\ +16{,}10]\) \\
Controle sintético  & \(-34{,}64\) & 16{,}73 & \([-67{,}43;\ -1{,}85]\) \\
DiD                 & \(-21{,}85\) & 19{,}23 & \([-59{,}54;\ +15{,}84]\) \\
\bottomrule
\end{tabular}
\end{table}

O controle sintético puro é o único estimador deste trabalho cujo intervalo
exclui o zero, e sua magnitude supera --- por pouco --- o efeito mínimo
detectável do desenho. \textbf{Ele não é reportado como achado}, por três
razões, e registrá-las é o ponto.

\begin{enumerate}
  \item Os três estimadores rodam sobre a \textbf{mesma matriz}. Divergirem
    quanto a excluir o zero não é evidência a favor do que exclui --- é
    evidência de que o resultado depende do estimador. A amplitude entre eles é
    de 14,5 g.
  \item Nenhum dos três é confirmatório. A pré-especificação declara como
    confirmatórios a curva dose-resposta e o óbito fetal, sob correção de Holm;
    tudo nesta subseção é exploratório declarado.
  \item Selecionar, entre muitos procedimentos, aquele que cruzou a linha é
    exatamente a prática que a pré-especificação foi escrita para impedir. O
    documento existe para tornar essa escolha visível, e a torna.
\end{enumerate}

\subsection{O que este ensaio pode e não pode afirmar}
\label{sub:balanco}

O quadro é coerente entre estimadores: \(-36\) g pela curva, \(-16\) a \(-33\) g
pelo \emph{event study} binário, \(-20\) g pelo controle sintético com DiD,
\(-22\) g pelo DiD simples. Todos negativos, quase todos cobrindo o zero, nenhum
superando o efeito mínimo detectável do desenho com margem.

\begin{quote}
\textbf{O que se pode afirmar:} o desenho é compatível com um efeito do
banimento sobre o peso ao nascer de magnitude até cerca de 60 g em módulo, e
não distingue esse efeito do zero. É limite superior informativo, e é assim que
deve ser lido.
\end{quote}

\paragraph{E há uma razão substantiva para o sinal negativo não ser
surpreendente.}
A seleção para nascimento vivo predizia exatamente isso. Se o banimento reduz o
óbito fetal, fetos marginais que antes não nasciam passam a nascer, e entram na
cauda inferior da distribuição de peso --- o efeito sobre o peso médio vem
atenuado ou invertido \citep{sorensen2025birth}. Foi por essa razão que o óbito
fetal entrou como segundo confirmatório, e não como acessório. O coeficiente
correspondente é de \(+0{,}0002\), com intervalo cobrindo o zero com folga: não
sustenta a história de seleção, e também não a exclui, porque 98,7\% das células
de óbito fetal têm menos de cinco eventos.

\paragraph{E há uma leitura concorrente deste diagnóstico inteiro.}
A §\ref{sub:equipamento} mostra que apenas trinta e seis estabelecimentos no
estado usavam aeronave, e que \textbf{quinze dos dezessete} municípios do decil
superior da dose não tinham nenhuma. Se esse padrão de 2006 ainda valia em 2018
--- o que não se pode verificar, porque o Censo seguinte não publicou o corte ---,
então parte do que esta seção lê como ausência de poder é, na verdade,
\textbf{tratamento atribuído a quem não o recebeu}.

As duas leituras pedem coisas opostas, e é por isso que distingui-las importa
mais do que escolher entre elas. Falta de poder pede mais unidades tratadas ---
a via declarada no parágrafo seguinte. Erro de medida do tratamento pede
\emph{outra variável de dose}, e encorpar o grupo tratado \textbf{pioraria} o
problema, diluindo ainda mais.

\begin{quote}
\textbf{Esta é a pergunta mais valiosa que o trabalho deixa em aberto}, e ela não
existia antes de olhar o equipamento: \emph{a dose por área plantada mede o que
o banimento removeu?} O cadastro aeroagrícola estadual, registrado como pendente
na §\ref{sub:aquisicao}, é a fonte que a responderia --- e passa, com isso, de
conveniência a prioridade.
\end{quote}

\paragraph{O que faria o desenho falar.}
O caminho declarado não é trocar o desfecho nem descer de unidade geográfica ---
as duas coisas estão vedadas na pré-especificação. É encorpar o grupo tratado,
baixando o corte de dose para além do decil superior, e é a via que o roteiro já
ratificava antes de haver estimativa.

\paragraph{E a âncora da literatura aponta um segundo caminho, que não estava
declarado.}
O achado central de \citet{reynier2025glyphosate} não é a média: é que o efeito
se concentra \textbf{doze vezes} mais no decil inferior de peso esperado ---
75 g de perda ali contra 6 g no decil superior. Posta essa distribuição contra o
efeito mínimo detectável deste desenho, aparece uma assimetria que o diagnóstico
agregado esconde:

\begin{table}[htbp]
\centering
\caption{Magnitudes da âncora contra o MDE deste desenho (33,4 g)}
\label{tab:mde-cauda}
\small
\begin{tabular}{lrc}
\toprule
\textbf{Alvo em \citet{reynier2025glyphosate}} & \textbf{Magnitude} & \textbf{Detectável aqui?} \\
\midrule
Efeito médio                       & 23--32 g   & não \\
Decil inferior de peso esperado    & 75 g       & sim \\
Exposição no percentil 90          & 146--243 g & sim \\
\bottomrule
\end{tabular}
\end{table}

Se o mesmo padrão distributivo valer no Ceará, \textbf{este desenho tem poder
para a cauda vulnerável e não para a média} --- e o problema deixa de ser de
tamanho de amostra para ser de escolha de estimando.

Isso \textbf{não} é reportado como resultado, e nem sequer como análise feita.
Não foi declarado antes do dado, e promovê-lo agora seria precisamente a prática
que a pré-especificação existe para impedir. Fica registrado como hipótese a
pré-especificar no ciclo seguinte, o que é o único estatuto honesto que ela pode
ter neste texto.
